% Transcription — fonds Alexandre Grothendieck, shelfmark 136, batch 5 (pages 81-100).
% Established from the facsimile archives/batches/136/batch-05.pdf (cover sheet
% excluded; batch page 1 = archivists' page 81), per the skill
% transcribe-grothendieck.
%
% Pass: Opus 5.5 (claude-opus-5-5), 2026-10-03 — first pass, unchecked against the pages by a human.
%
% Pages 83, 85, 87, 89, 91, 92, 94, 96, 98 and 100 carry no writing of his
% (typescripts by other hands, or blank) and are absent. Pages 81, 82, 84, 86,
% 88, 90, 93, 95, 97 and 99 are his manuscript.
%
% Contents, in outline:
%   Pages 81-86 — the bicomplex K^{..} attached to an immersion u : L -> Phi of
%                 a flat k-Module: functoriality in u, the list a)-g) of its
%                 structures (bigraded algebra, divided powers, differential,
%                 K^{..} = Gamma^* Phi (x) Lambda^* Psi), the data (k, L, Phi, u)
%                 and the homomorphisms between two such; the question of two
%                 homomorphisms lambda, mu.
%   Pages 88-90 — X a topos, k a Ring: the lemma on a « universal » mono
%                 u_N : N -> \check N; the complex DR_pd^{**}(k,u); its
%                 independence of the choice of S -> N (transitivity); the
%                 « condition de De Rham » on a Ring k; algebras A^{**} in
%                 DR_pd^{**}(X).
%   Pages 93-99 — category of special De Rham algebras (with divided powers):
%                 definition, the categories DRsp(X,k)_0, DRspfl, DRspufl; the
%                 universal-flasque objects DR(Phi,T); the question whether
%                 flasque + flat implies homotopic to 0, the projective topos P
%                 and the Godement-type resolution u_n^*(K) = K_n = K(Delta_n).
%
% Notation fixed for the batch: \Phi, \Psi for his capital phi, psi; L for the
% flat Module (on p. 81 the letter is drawn with a doubled stroke and looks like
% a Z: it is the L of p. 84); K^{\cdot\cdot} for the bicomplex; DR_{pd} for his
% De Rham « à puissances divisées »; \mathbb{DR} for his double-struck DR
% (a category); \mathcal{DR} for his script DR (a sheaf); k\{T\} for the
% divided-power polynomial algebra; \widehat{\Delta} for the topos of
% simplicial sets. His « ½ » prefix (= semi-) is kept as written.
%
\documentclass[11pt,a4paper]{article}
\input{../preamble/grothendieck}

\folder{136}
\batch{5}
\pages{81}{100}
\foldertitle{Complexe de De Rham à puissance divisée [conférence de 1976 à l’IHÉS] : notes manuscrites (s.d.).}
\dating{[à partir de 1975-1976]}
\watermark{Édition de démonstration}

\begin{document}

\page{81}
\note{la phrase commence sur la page précédente, hors de ce lot}
est un isomorphisme \add{en cohomologie} (en bidegré $(i,j)$ avec $j\leq i$),
tandis que la coh.\ du premier membre est nulle pour $j>i$.

On peut l'exprimer comme un isomorphisme du \ill{} avec un certain \ill{} du
second…

\textbf{Fonctorialités.} a) en $k$, clair.

b) \emph{En $u$} \struck{\ill{}}. Soit
\[
  u'\colon L\longrightarrow\Phi'
\]
une deuxième immersion\,; s'il existe un hom.\ $\lambda\colon\Phi\to\Phi'$
rendant commutatif

\begin{tikzcd}
 & L \arrow[dl, "u"'] \arrow[dr, "{u'}"] & \\
\Phi \arrow[rr, "\lambda"'] & & {\Phi'}
\end{tikzcd}

i.e.\ $\lambda u=u'$, on en conclut $\bar\lambda\colon K^{\cdot\cdot}_{u}\to
K^{\cdot\cdot}_{u'}$ et $\bar\lambda_{h}\colon K^{\cdot\cdot}_{u,h}\to
K^{\cdot\cdot}_{u',h}$, compatible avec les structures, qui est un isom.\ dans
telle catégorie dérivée, et le reste après application de $\Gamma(X,-)$ ou de
$f_{*}$.

Supposons qu'on ait un autre, soit $\mu$, tel que $\mu u=u'$, donc, posant
$\mu=\lambda+\Theta$, on trouve $\Theta u=0$
\note{la page s'arrête ici}

\page{82}
$K^{\cdot\cdot}$ \quad Algèbre de $\mathrm{DR}$ à p.d.\ élémentaire

a) Anneau bigradué, \struck{\ill{}} alterné pour \add{« deg.\ ext »}
\marginal{$\surd$}

b) \struck{\ill{}} à bidegré $\geq0$

\quad Structure à p.d.\ sur l'idéal $\sum\limits_{i+j>0}K^{ij}$

c) différentielle d'Anneau gradué (pour le grad.\ ext) en changeant pour le
degré total, etc.\ comp.\ avec p.d.

d) $d\colon K^{1,0}=\Phi\longrightarrow K^{11}=\Psi$ est un épi.

e) Posant $k=K^{00}$, $K^{**}$ est l'algèbre \add{alternée à p.d.}\ libre
engendrée par $\Phi$ et $\Psi$, i.e.\ $K^{**}\simeq\Gamma^{*}\Phi\otimes
\bigwedge^{*}\Psi$ [Il en résulte que $K^{ij}=0$ si $j>i$]

f) \struck{\uncertain{Mono-plat}} $\Phi$ et $\Psi$ sont $k$-plats

(donc aussi $L=\operatorname{Ker}d=\mathcal{H}^{1\,0}(K^{\cdot\cdot})$)

Il résulte déjà de ceci que
\[
  H^{*i}(K^{\cdot\cdot})=
  \begin{cases}
    \Gamma^{*}L & \text{si } i=0\\
    0 & \text{si } i>0
  \end{cases}
\]
donc que $K^{\cdot\cdot}$ est une résolution de $\Gamma^{*}L$. On va supposer
le plus souvent

g) $L$ est un $k$-Module \struck{inversible} \add{loc.\ libre} de rang $1$.

(voire libre de rang $1$ et trivialisé…)

\page{84}
\struck{Soient $K^{\cdot\cdot}$ et $K'^{\cdot\cdot}$}

Finalement, la donnée de $K^{\cdot\cdot}$ équivaut à celle de

$\alpha$) $k$, Anneau commutatif sur $X$

$\beta$) $L$, $k$-Module plat (voire inversible, voire $k$ lui-même…)

$\gamma$) Une immersion de $L$ dans un $k$-module $\Phi$, $u\colon L\to\Phi$.

Soit de même $K'^{\cdot\cdot}$ défini par $(k',L',\Phi',u')$. La donnée d'un
hom.\ $\lambda$ de $K^{\cdot\cdot}$ dans $K'^{\cdot\cdot}$, compatible avec les
structures, équivaut à la donnée d'un système d'hom.
\[
  k\xrightarrow{\ \lambda_{0}\ }k'\qquad\text{hom.\ d'Anneaux}
\]
\struck{$K^{\cdot\cdot}\to K'^{\cdot\cdot}$}
\[
  \Phi\xrightarrow{\ \lambda_{\Phi}\ }\Phi'
\]
hom.\ $\lambda_{0}$-½linéaire
avec $\lambda_{\Phi}(u(L))\subset u'(L')$, \struck{\ill{}} donc
$\lambda_{\Phi}$ induit un hom.\ ($\lambda_{0}$-½linéaire)
\note{« ½linéaire » : semi-linéaire, au-dessus de $\lambda_{0}$}

\begin{tikzcd}
L \arrow[r, "\lambda_{L}"] \arrow[d] & {L'} \arrow[d] \\
\Phi \arrow[r] & {\Phi'}
\end{tikzcd}

\note{un premier tracé du même carré, à gauche, est biffé}

Dans le cas $L\simeq k$, $L'\simeq k'$ donnés, on peut oublier $L,L'$ et
retenir $\Phi,\Phi'$ avec sections marquées, et la condition sur
$\lambda_{\Phi}$

\page{86}
devient donc de transformer section marquée en section marquée.

Supposons qu'on ait deux hom.\ $\lambda,\mu$, correspondant\,: un
$k\xrightarrow{\lambda_{0}=\mu_{0}}k'$ et deux hom.
\[
  \Phi\overset{\lambda_{\Phi},\,\mu_{\Phi}}{\rightrightarrows}\Phi'
\]
(ayant même effet sur la section marquée). Que peut-on dire des hom.\
$\lambda,\mu$ de $K^{\cdot\cdot}$ dans $K'^{\cdot\cdot}$\,?

\page{88}
$X$, topos

$k$ Anneau sur $X$

$k\overset{u}{\hookrightarrow}\Phi$ mono dans $k$-Modules, avec
\[
  \begin{cases}
    \text{(a)}\ \Psi=\operatorname{Coker}u\ \ k\text{-plat}\\
    \text{(b)}\ \forall\ k\text{-modules } N, M,\ \text{et } i>0,\\
    \qquad\Gamma^{i}(N\otimes_{k}\Phi)\otimes M\ \text{est flasque}
  \end{cases}
\]

\textbf{Lemme.} $\forall$ $k$-module $N$, $\exists$ \uncertain{homom.}\
$N\xrightarrow{u_{N}}\check N$, qui est un $k$-mono « universel » (vérifiant
mêmes \uncertain{propriétés}\ \ill{}) \struck{\ill{}}\ \uncertain{satisfaisant}\
\struck{$N$}\,: (a) \ill{} \uncertain{plat} $\Rightarrow$
$\operatorname{Coker}u_{N}$ plat, et (b$'$) $\forall i>0$, et $k$-modules
\struck{$N$}, $M$
\[
  \Gamma^{i}_{k}\bigl(\check N\,\text{\struck{$\otimes_{k}\Phi$}}\bigr)\otimes_{k}M
  \quad\text{flasque}
\]
On prend $\check N=\Phi\otimes_{k}N$, $u_{N}=u\otimes\mathrm{id}_{N}$, donc (a)
résulte de (a), et (b$'$) de (b).

$\mathcal{DR}^{**}_{pd}($\struck{\ill{}}\,$k,u)=\Gamma^{*}(\Phi)\otimes_{k}\bigwedge^{*}\Psi$
Algèbre bigraduée (degré complet $+$ alternée pour deg.\ ext), augm.\ cat.\
degré total) (différentielle à p.d.\ sur idéal des bi-degrés $>0$…) ---
résolution de $k\{T\}$ (où deg ext $T=0$, deg tot $T=1$)
\note{lecture des indices de degré incertaine}

\[
  \mathrm{DR}^{**}_{pd}(k,u)=\Gamma\bigl(X,\mathcal{DR}^{**}_{pd}(k,u)\bigr).
\]

Soit $\mathbb{DR}^{**}$ la cat.\ des \add{anneaux} bigradués (à deg
$\geq0$) \add{alternés, diff.\ à p.d.}\ \struck{anticommutatifs, engendrés
par bidegré $(1,0)$}, \add{avec} \struck{\ill{}} \add{\uncertain{élément}\ de
degré $(1,0)$} \ill{}…\ On a défini un système transitif d'isom.\ dans
$\mathbb{DR}^{**}$ entre les $\mathrm{DR}^{**}_{pd}(k,u)$, pour $k$ fixé et
$u$ variable.

N.B.\ hom.\ pour \ill{}

\begin{tikzcd}[column sep=small, row sep=small]
 & \Phi \arrow[d, "\lambda"] \\
k \arrow[ur, "u"] \arrow[r, "{u'}"'] & {\Phi'}
\end{tikzcd}

d'où $\lambda^{**}\colon\mathrm{DR}^{**}_{pd}(k,u)\to\mathrm{DR}^{**}_{pd}(k,u')$,
donc $\lambda^{**}$ est un iso.\ dans $\mathbb{DR}^{**}$

\note{il s'agit d'un hom.\ $\lambda$ quelconque, non de l'iso.\ donné par le lemme\,; cf.\ ce qui suit}

\begin{tikzcd}[column sep=small, row sep=small]
 & {\Phi'} \arrow[dr, hook] & & \\
k \arrow[ur, hook, "u"] \arrow[dr, hook, "{u'}"'] & & S \arrow[r, hook] & N \\
 & \Phi \arrow[ur, hook] & &
\end{tikzcd}

($S=\Phi\amalg_{k}\Phi'$, et $S\hookrightarrow N$ donné par le lemme), d'où
$\mathrm{DR}(u)\xrightarrow{\ \sim\ }\mathrm{DR}(N)\xleftarrow{\ \sim\ }\mathrm{DR}(u')$,
d'où $\mathrm{DR}(u)\simeq\mathrm{DR}(u')$ dans $\mathbb{DR}^{**}$

indépendant du choix de $S\hookrightarrow N$ avec $N$ comme \ill{} dans le
lemme\,:

\begin{tikzcd}[column sep=small, row sep=small]
 & N \arrow[dr, hook] & & \\
S \arrow[ur, hook] \arrow[dr, hook] & & T \arrow[r, hook] & {N''} \\
 & {N'} \arrow[ur, hook] & &
\end{tikzcd}

($T=N\amalg_{S}N'$, $T\hookrightarrow N''$ donné par le lemme), d'où

\begin{tikzcd}[column sep=small, row sep=small]
\mathrm{DR}(u) \arrow[dr] & & \mathrm{DR}(N) \arrow[dr, "\sim"] & \\
 & \mathrm{DR}(S) \arrow[ur] \arrow[dr] & & \mathrm{DR}(N'') \\
\mathrm{DR}(u') \arrow[ur] & & \mathrm{DR}(N') \arrow[ur, "\sim"'] &
\end{tikzcd}

Transitivité évidente par diagramme

\page{90}
\drawing{à gauche, un schéma à traits sans flèches\,: $k$ relié par $u$, $u'$, $u''$ à $\Phi$, $\Phi'$ (écrit deux fois, un tracé biffé), $\Phi''$\,; $\Phi,\Phi'$ vers $N$, $\Phi',\Phi''$ vers $N'$, puis $N$ et $N'$ vers $N''$}

d'où

\begin{tikzcd}[column sep=small, row sep=small]
 & \Phi \arrow[dr] & \\
k \arrow[ur] \arrow[r] \arrow[dr] & {\Phi'} \arrow[r] & {N''} \\
 & {\Phi''} \arrow[ur] &
\end{tikzcd}

\[
  \begin{aligned}
    &\mathrm{DR}(\Phi)\ \cdots\\
    &\mathrm{DR}(\Phi'')\xrightarrow{\ \simeq\ }\mathrm{DR}(N'')\\
    &\mathrm{DR}(\Phi'')\nearrow\simeq
  \end{aligned}
\]
\note{les trois flèches vers $\mathrm{DR}(N'')$ sont esquissées\,; la deuxième ligne porte sans doute $\mathrm{DR}(\Phi')$, que la page écrit comme $\Phi''$}

Argument analogue pour $f_{*}\bigl(\mathcal{DR}^{**}_{pd}(k,u)\bigr)$, où
$f\colon X\to Y$ morphisme de topos.

\emph{Conditions de De Rham} \emph{sur Anneau $k$}\,: la condition
d'existence d'un $u\colon k\to\Phi$ tel que…

Si $k$ satisfait condition de De Rham, de même tt $k$-Algèbre. Donc le plus
beau, c'est que $\mathbb{Z}_{X}$ y satisfasse\,! C'est vrai pour

\quad a) le topos SS $\widehat{\Delta}$

\quad b) le topos défini par un espace topologique \uncertain{paracompact}.

\emph{Conjecture.} Est-ce vrai pour $\widehat{\Delta}^{\circ}$ (topos des
\uncertain{ens.}\ cosimpliciaux)\,? \emph{Non}
\note{« Non » souligné trois fois, sans doute ajouté après coup}

Soit $A^{**}\in\operatorname{Ob}\mathbb{DR}^{**}_{pd}(X)$ une Algèbre bigr.\
\uncertain{donc} à degrés positifs, alternée pour le degré ext, diff.\ à p.\
div.\ …\ Supposons que $A^{00}=k$ satisfasse la condition de De Rham, soit
$u\colon k\to\Phi$, d'où $\mathcal{DR}^{**}_{pd}(k,u)$. Considérons alors
\[
  \overline{A^{**}}=\mathcal{DR}^{**}_{pd}(A^{**})
  \overset{\mathrm{déf}}{=}\mathcal{DR}^{**}_{pd}(k,u)\otimes_{k\{T\}}A^{**},
\]
l'augmentation $k\{T\}\to\mathcal{DR}^{**}_{pd}(k,u)$ (qui est un
quasi-iso.\ \uncertain{universel}) définit un hom.
\[
  A^{**}\longrightarrow\overline{A^{**}}
\]
qui est un quasi-iso.\ \add{\uncertain{Ensuite}} \struck{\ill{}}
\note{la phrase s'arrête au bas de la page\,; un mot souligné en fin de ligne, peut-être « Ensuite » écrit sur un mot biffé}

\page{93}
\section{Cat.\ des Alg.\ de De Rham spéciales}
\note{titre de sa main, encadré en haut à droite de la page}

\struck{(à p.\ div.\ \ill{})}

\emph{Structure De Rham} (à p.d.\ -- sous-entendu)\,:

\quad Anneau bigradué [à degrés -- total et ext -- $\geq0$] alterné pour degré
ext
\[
  A^{ij}\neq0\ \text{\struck{$\Leftarrow$}}\Rightarrow 0\leq j\leq i
\]
\quad degré complémentaire $=i-j\geq0$

différentiel [différentielle de bidegré $0,1$\,; c'est une diff.\ pour le
degré ext\,:
\[
  d(xy)=(dx)y+(-1)^{\deg\mathrm{ext}\,x}x\,dy\ ]
\]
à p.\ div.\ sur l'idéal
\[
  \sum_{i+j>0}A^{ij}=A^{+}
\]
[p.d.\ \struck{\ill{}} compatible avec 2 grad.\ ($\gamma^{d}\colon A^{ij}\to
A^{di,dj}$) et la différentielle ($d\,x^{(i)}=x^{(i-1)}dx$)]

élément marqué $T\in Z^{10}(A^{**})\ \bigl(\xrightarrow{\ \simeq\ }H^{10}(A^{**})\bigr)$]

[Cohomologie\,: $H^{*i}(A^{**})\simeq[h^{i}\otimes_{k}k\{T\},0]$, $k=A^{00}$]
\note{lecture de cette ligne incertaine}

Structure de De Rham \emph{spéciale}\,:
\[
  \begin{cases}
    k\{T\}\xrightarrow{\ \sim\ }H^{*0}(A^{**})\\
    H^{*i}(A^{**})=0\quad(i>0)\\
    \text{les } A^{ij}\text{ sont $k$-plats}
  \end{cases}
\]
[NB c'est utile pour la stabilité par $\otimes_{k}k'$]
\note{la condition $H^{*i}=0$ est écrite sans précision sur $i$\,; « $(i>0)$ » est notre lecture du contexte}

\marginal{\struck{\ill{}} Univ.\ flasque\,: $\to$ flasques. \struck{\ill{}}}

[Structure de De Rham flasque\,: les $A^{ij}$ ($1\leq i$) flasques.]
\note{ligne ajoutée en interligne\,; lecture incertaine}

On parle de $k$-structure de De Rham si $k=A^{00}$ est donné.

\struck{\ill{}} $\mathfrak{X}$ topos. \struck{\ill{}}

$\mathbb{DR}\mathrm{sp}(\mathfrak{X},k)_{0}$\,: structures de De Rham
spéciales sur $\mathfrak{X}$ -- avec comme flèches celles qui respectent
\struck{toutes} les structures (y compris les structures)

$\mathbb{DR}\mathrm{spfl}(\mathfrak{X},k)_{0}$\,: « \quad spéciales et flasques

$\mathbb{DR}\mathrm{spufl}(\mathfrak{X},k)_{0}$\,: « \quad --- \quad univ.\ flasques

$\mathbb{DR}\mathrm{sp}(\mathfrak{X},k)$, $\mathbb{DR}\mathrm{spfl}(\mathfrak{X},k)$,
$\mathbb{DR}\mathrm{spufl}(\mathfrak{X},k)$ \add{\uncertain{de fractions}}\,:
catégories \uncertain{déduites} en rendant inversibles les flèches
\uncertain{qui sont} des quasi-isom.\ --- NB on \ill{}
\note{la phrase se poursuit au haut de la page suivante, p.~95}

\page{95}
\note{tableau des $A^{ij}$ en haut de la page\,; les flèches verticales sont les différentielles}

\begin{tikzcd}[column sep=small, row sep=small]
0 & 0 & A^{22} & \cdots \\
0 & A^{11} & A^{21} \arrow[u] & \cdots \\
k & A^{10} \arrow[u] & A^{20} \arrow[u] & \cdots
\end{tikzcd}

Soit $A^{**}\in\operatorname{Ob}\mathbb{DR}\mathrm{sp}(\mathfrak{X},k)$\,;
$T\in\Gamma_{\mathfrak{X}}(A^{10})$ définit $k\to A^{10}=\Phi$, $\lambda\mapsto\lambda T$, avec suite exacte

\begin{tikzcd}[column sep=small, row sep=small]
0 \arrow[r] & k \arrow[r] & \Phi \arrow[r] \arrow[d, no head, "\Vert" description] & \Psi \arrow[r] \arrow[d] & 0 \\
 & & A^{10} & A^{20} &
\end{tikzcd}

\note{la flèche verticale sous $\Psi$ va vers un terme que nous lisons $A^{20}$, peut-être $A^{11}$}

Comme $\Psi$ est $k$-plat, $\Phi$ \ill{}\ \ill{}, et on a \ill{}
\[
  \mathrm{DR}(\Phi,T)=\underbrace{\Gamma^{*}(\Phi)}_{\text{deg compl}}
  \otimes_{k}\underbrace{\textstyle\bigwedge^{*}\Psi}_{\text{deg ext}}
\]
\note{au-dessus de « $\mathrm{DR}$ » un mot ajouté, \ill{}}

\ill{}\ \ill{} est \ill{} et on a un $k$-\ill{}\ \ill{}

\struck{donc la cat.\ de}

Un hom.\ $A^{**}\to A'^{**}$ \ill{}

\begin{tikzcd}
\mathrm{DR}(\Phi,T) \arrow[r] \arrow[d] & A^{**} \arrow[d] \\
\mathrm{DR}(\Phi',T') \arrow[r] & {A'^{**}}
\end{tikzcd}

Cela implique \struck{\ill{}} ainsi que la cat.\ \add{dérivée}
$\mathbb{DR}\mathrm{sp}(\mathfrak{X},k)_{0}$ est équivalente canoniquement\,:
la \struck{\ill{}} cat.\ qui se déduit des \add{\uncertain{dérivées}} de la forme
$\mathrm{DR}(\Phi,T)$, qui est aussi celle des \ill{} exactes
$0\to k\to\Phi\to\Psi\to0$ avec $\Psi$ plat sur $k$.

\page{97}
On sait que cette dernière est équivalente à la cat.\ ponctuelle.
Donc $\mathbb{DR}\mathrm{sp}(\mathfrak{X},k)_{0}$ est équivalente à la cat.\
ponctuelle.

Si maintenant on travaille dans
\struck{(cat.\ $\mathbb{DR}\mathrm{spufl}$)} $\mathbb{DR}\mathrm{spufl}$ i.e.\
avec des $A^{**}$ flasques, \add{u.} $(A^{ij}$ flasque si $0<j)$ \struck{\ill{}}
\ill{}, alors $\Phi$ est flasque.
Supposons aussi que (\uncertain{les}) \ill{} les $\Gamma^{i}_{k}\Phi$ ($i>0$)
sont univ.\ flasques, alors $\mathrm{DR}(\Phi,T)$ est aussi flasque,
\struck{\ill{}} et on voit comme dessus que la cat.\ \uncertain{dérivée} est
équivalente à celle construite \ill{} \struck{la cat.\ des}
$\mathrm{DR}(\Phi,T)$, avec $\Phi$ univ.\ flasque et $\Phi/k\cdot T$ plat,
i.e.\ avec les suites exactes $0\to k\to\Phi\to\Psi\to0$, \ill{}
$\Phi$ univ.\ flasque et $\Psi$ $k$-plat\,; or cette cat.\ dérivée est
\uncertain{encore} la catégorie ponctuelle. On a donc\,:

\emph{Prop.} La cat.\ $\mathbb{DR}\mathrm{sp}(\mathfrak{X},k)_{0}$ est
équivalente à la cat.\ ponctuelle. Il en est de même de la cat.\
$\mathbb{DR}\mathrm{spufl}(\mathfrak{X},k)_{0}$, \add{$k$-\uncertain{modules}}
\struck{\ill{}} si on suppose que pour \ill{} \struck{\ill{}} $\Phi$ plat sur
$k$ univ.\ flasque, les $\Gamma^{i}\Phi$ ($i>0$) sont aussi univ.\ flasques.

\page{99}
Cette dernière condition est satisfaite pour le topos \add{$(\mathrm{Ss})=$}
$\widehat{\Delta}$, car pour celui-ci
\[
  \text{flasque}+\text{plat}\Longrightarrow\text{homotope à $0$, donc}
\]
\[
  \Downarrow
\]
(\uncertain{univ.\ projectif})
univ.\ flasque et les $\Gamma^{i}(\Phi)$ univ.\ flasques ($i>0$).
\marginal{faux tel quel, à arranger\,!}
\note{l'accolade sous « flasque $+$ plat » renvoie à « univ.\ projectif », souligné\,; la note marginale porte sur cette implication}

Soit $P$ le topos ponctuel, et
\[
  i_{n}\colon P\longrightarrow\widehat{\Delta}(\mathrm{S}_{s})
\]
le morphisme de topos associé au foncteur
\[
  u_{n}^{*}(K)=K_{n}=K(\Delta_{n}).
\]
Soit $k$ un anneau (commutatif, \uncertain{disons}).

\emph{Lemme.} $\forall$ $k$-module $M$, on a un
\[
  i_{0*}(M)\simeq\Phi_{*}\otimes_{\mathbb{Z}}M\simeq\Phi_{*k}\otimes_{k}M
\]
où $\Phi_{*}\in\operatorname{Ob}\bigl(\widehat{\Delta}_{ab}\simeq
\underline{\mathrm{Hom}}(\Delta^{\circ},\mathrm{Ab})\bigr)$ est défini par
$\Phi_{*}(I)=\mathbb{Z}^{I}$, donc $\Phi_{*k}(I)\simeq k^{I}$.

En effet, on a par construction de l'image directe
\[
  i_{0*}(M)(I)\simeq M(i_{0}^{*}(I))\simeq M^{I}
\]
\[
  \mathrm{Hom}(I,i_{0*}(M))=\mathrm{Hom}\bigl(\underbrace{i_{0}^{*}(I)}_{\mathrm{Hom}(\Delta_{0},I)\simeq I},M\bigr)
\]
OK.

\struck{\ill{}} \emph{Corollaire.} Pour tout $k_{*}$-Module $A_{*}$,
\[
  \mathrm{Hom}(A_{*},\Phi_{*}\otimes_{k}M)\simeq\mathrm{Hom}_{k}(A_{0},M)
\]
en particulier on a un hom.\ can.\ (pour $M=A_{0}$)
\[
  \text{\struck{$\mathrm{Hom}$}}\ A_{*}\longrightarrow\Phi_{*}\otimes_{\mathbb{Z}}A_{0}
  \simeq\Phi_{*k}\otimes_{k}A_{0}\quad(\simeq i_{0*}i_{0}^{*}A_{*})
\]
Appliquons ceci au cas d'une immersion
\[
  k_{*}\overset{u}{\hookrightarrow}A_{*}=A^{1,0}_{*}
\]
provenant d'une $A^{**}_{*}\in\operatorname{Ob}\mathbb{DR}\mathrm{sp}(\mathfrak{X}_{s},k_{*})$
\note{la phrase se poursuit au-delà de ce lot}

\end{document}
